\documentclass[11pt,a4paper]{article}
\usepackage[margin=1in]{geometry}
\usepackage{times}
\usepackage{latexsym}
\usepackage{amsmath,amssymb}
\usepackage[colorlinks=true,citecolor=blue,urlcolor=blue]{hyperref}
\usepackage{natbib}
\renewcommand{\bibsection}{\section*{References}}
\usepackage{booktabs}
\usepackage{multirow}
\usepackage{graphicx}
\usepackage{xcolor}
\usepackage{algorithm}
\usepackage{algorithmic}

\title{Cardinality-Constrained MCKP for Per-Head Mixed-Precision KV Cache Quantization}

\author{Anonymous}

\begin{document}
\maketitle

%======================================================================
\begin{abstract}
KV cache quantization is a key technique for reducing memory footprint during large language model inference.
Existing approaches apply uniform bit-widths across all attention heads or rely on continuous relaxations that ignore the discrete nature of hardware-supported precisions.
We show that per-head quantization sensitivity in modern LLMs is extremely non-uniform: Layer~0 heads exhibit $100\times$ higher profiling delta than average, and fewer than 4\% of heads account for over 70\% of total perplexity degradation under uniform quantization.
We formulate per-head bit-width allocation as a \emph{cardinality-constrained multiple-choice knapsack problem} (MCKP) over the discrete action set $\{2,4,8,16\}$-bit.
The cardinality constraint---limiting the number of heads assigned extreme quantization---acts as principled regularization against compounding effects that independent per-head profiling cannot capture.
On Qwen2.5-7B at $4\times$ KV cache compression, constrained MCKP achieves perplexity 11.08, a 28.3\% improvement over uniform 4-bit (15.45), while unconstrained optimization catastrophically degrades to 36.92.
Results generalize across Llama-3.1-8B (24.2\% improvement) and Mistral-7B-v0.3 (22.2\% improvement), and across calibration datasets (rank correlation 0.94 between WikiText-2 and C4 sensitivity profiles).
The solver automatically discovers architectural sensitivity structure, assigning 8-bit protection to all Layer~0 heads without any domain-specific heuristic.
\end{abstract}

%======================================================================
\section{Introduction}

Large language models (LLMs) with billions of parameters increasingly rely on KV caching for efficient autoregressive generation.
At long context lengths, the KV cache dominates GPU memory, motivating quantization to lower bit-widths.
The standard approach applies a uniform precision---typically 4-bit---to all KV heads.
However, not all heads contribute equally to model quality: some are highly sensitive to quantization noise while others tolerate extreme compression with negligible impact.

We make three observations that motivate our approach.
\textbf{First}, per-head sensitivity is remarkably heterogeneous.
Profiling Qwen2.5-7B reveals that Layer~0 heads have quantization deltas $100\times$ larger than the average, and just 3.6\% of heads (4 out of 112) account for over 70\% of the total perplexity degradation under uniform 4-bit quantization.
This pattern holds across architectures: Llama-3.1-8B and Mistral-7B-v0.3 exhibit the same Layer~0 dominance despite different GQA ratios and attention mechanisms.

\textbf{Second}, naive exploitation of this heterogeneity fails catastrophically.
An unconstrained solver that trusts per-head independent profiling and freely assigns 2-bit to ``insensitive'' heads produces perplexity 36.92---$3.3\times$ worse than uniform 4-bit.
The root cause is \emph{compounding}: when many heads are simultaneously assigned extreme quantization, per-head errors compound through the residual stream in a strongly non-additive manner.

\textbf{Third}, the compounding effect exhibits a sharp threshold.
Sweeping the number of 2-bit heads from 0 to 112 on Qwen2.5-7B reveals an elbow at $k \approx 8$--$12$: below this threshold, each additional 2-bit head costs only 0.055 PPL; above it, the marginal cost jumps to 1.29 PPL---a $23\times$ slope increase (SNR $= 7.3$).

These observations lead to our formulation: per-head bit-width allocation as a \emph{cardinality-constrained multiple-choice knapsack problem} (MCKP).
Each KV head is an item with four discrete bit-width options $\{2, 4, 8, 16\}$, a knapsack budget constraining total bit cost to match the target compression ratio, and a cardinality constraint limiting the number of heads at extreme (2-bit) quantization.
The cardinality constraint regularizes the profiler's systematic overconfidence in extreme assignments, converting a brittle unconstrained optimizer into a robust allocation mechanism.

Our constrained MCKP solver achieves 11.08 perplexity on Qwen2.5-7B at $4\times$ compression---28.3\% better than uniform 4-bit and within $<$1\% of a brute-force oracle.
Results are consistent across three model families, two calibration datasets, three compression ratios, and downstream tasks (MMLU, LongBench).

%======================================================================
\section{Related Work}

\paragraph{KV Cache Quantization.}
KIVI~\cite{kivi2024} and KVQuant~\cite{kvquant2024} apply per-channel or per-token quantization at uniform bit-widths.
ARKV~\cite{arkv2026} introduces per-layer heuristic precision selection but lacks per-head granularity and optimality guarantees.

\paragraph{Rate--Distortion Approaches.}
RateQuant~\cite{ratequant2026} formulates per-head allocation via reverse waterfilling, achieving R-D optimal solutions under continuous bit-width and additive distortion assumptions.
However, hardware supports only discrete precisions ($\{2,4,8,16\}$-bit), and rounding continuous solutions introduces suboptimality.
More critically, the additive distortion model breaks when many heads are simultaneously assigned extreme quantization due to compounding effects.
Our work addresses both limitations: discrete MCKP operates natively in the hardware action space, and cardinality constraints regularize compounding non-additivity.

\paragraph{Mixed-Precision Quantization.}
Mixed-precision methods for weight quantization~\cite{haq2019,hawq2020} solve similar allocation problems but target model weights rather than the dynamically generated KV cache.
The key difference is that KV cache quantization is an \emph{inference-time} decision that must be made once and deployed, making solver efficiency and calibration robustness critical.

%======================================================================
\section{Method}

\subsection{Per-Head Distortion Profiling}

Given a pretrained model with $N$ KV heads and a calibration corpus, we estimate the sensitivity of each head $i$ to quantization.
For each head $i$ and candidate bit-width $b \in \mathcal{B} = \{2, 4, 8, 16\}$, we measure the \emph{profiling delta} $\delta_{i,b}$: the change in model perplexity when head $i$ alone is quantized to $b$-bit while all other heads remain at full precision.

Formally, let $\text{PPL}(\mathbf{b})$ denote the perplexity under bit-width assignment $\mathbf{b} = (b_1, \ldots, b_N)$, and $\mathbf{b}^{16}$ denote the all-16-bit (FP16) assignment.
The profiling delta is:
\begin{equation}
\delta_{i,b} = \text{PPL}(\mathbf{b}^{16}_{-i}, b_i = b) - \text{PPL}(\mathbf{b}^{16})
\end{equation}
where $\mathbf{b}^{16}_{-i}$ keeps all heads except $i$ at 16-bit.

This per-head independent profiling is computationally efficient ($4N$ forward passes on a small calibration set) but assumes distortion additivity---an assumption we explicitly regularize in the solver.

\subsection{Constrained MCKP Formulation}

We formulate bit-width allocation as a constrained MCKP.
Each head $i \in \{1, \ldots, N\}$ must be assigned exactly one bit-width $b_i \in \mathcal{B}$.
Let $c_{i,b} = b / 16$ denote the normalized cost (fraction of FP16 memory) and $d_{i,b} = \delta_{i,b}$ the distortion.

\begin{align}
\min_{\mathbf{b}} &\quad \sum_{i=1}^{N} d_{i, b_i} \label{eq:obj} \\
\text{s.t.} &\quad \sum_{i=1}^{N} c_{i, b_i} \leq B \quad \text{(budget)} \label{eq:budget} \\
&\quad \sum_{i=1}^{N} \mathbf{1}[b_i = 2] \leq k \quad \text{(cardinality)} \label{eq:card} \\
&\quad b_i \in \mathcal{B} \quad \forall i \label{eq:discrete}
\end{align}

The budget constraint~\eqref{eq:budget} enforces the target compression ratio ($B = 0.25$ for $4\times$).
The cardinality constraint~\eqref{eq:card} limits the number of heads at 2-bit to at most $k$, regularizing compounding effects.

\subsection{Determining the Cardinality Bound $k$}

The cardinality bound $k$ is not a free hyperparameter but is determined empirically from the compounding curve.
We sweep $k$ from 0 to $N$ (or a representative subset), evaluate the resulting allocation's perplexity, and identify the \emph{elbow point} where the PPL-vs-$k$ curve transitions from approximately linear (additive regime) to superlinear (compounding regime).

On Qwen2.5-7B, the elbow occurs at $k \approx 12$ with SNR $= 7.3$, and the optimal allocation uses $k = 8$ (Section~\ref{sec:compounding}).
The sweep requires only $O(10)$ solver evaluations, each taking $< 0.1$s, plus the corresponding perplexity evaluations.

\subsection{Solver}

The constrained MCKP with budget and cardinality constraints is solved via LP relaxation with rounding, leveraging the problem's favorable structure: items have clear delta separation, the budget constraint is linear, and each head has only 4 options.
Wall-clock solver time is $< 0.05$s for all models tested, including Qwen2.5-72B with 560 KV heads.

%======================================================================
\section{Experiments}

\subsection{Setup}

\paragraph{Models.} Qwen2.5-7B (GQA 7:1, 4 KV heads/layer $\times$ 28 layers $= 112$ heads), Llama-3.1-8B (GQA 4:1, 8 KV heads/layer $\times$ 32 layers $= 256$ heads), and Mistral-7B-v0.3 (GQA 4:1, 8 KV heads/layer $\times$ 32 layers $= 256$ heads, sliding window attention).

\paragraph{Datasets.} WikiText-2 for primary evaluation and calibration; C4 for cross-dataset robustness analysis.

\paragraph{Baselines.} FP16 (unquantized), uniform $b$-bit, unconstrained MCKP (no cardinality constraint), and RateQuant~\cite{ratequant2026} (reverse waterfilling with rounding to $\{2,4,8,16\}$).

\paragraph{Profiling.} Default: 16 sequences $\times$ 512 tokens (8K tokens).
High-resolution variant: 64 sequences $\times$ 2048 tokens (128K tokens).

\subsection{Main Results: Cross-Model Consistency}
\label{sec:main}

Table~\ref{tab:main} presents results across three model families at $4\times$ KV cache compression.
Constrained MCKP consistently and substantially outperforms uniform 4-bit quantization: 28.3\% on Qwen2.5-7B, 24.2\% on Llama-3.1-8B, and 22.2\% on Mistral-7B-v0.3.

\begin{table}[t]
\centering
\small
\caption{Main results at $4\times$ KV cache compression on WikiText-2. PPL improvement is relative to uniform 4-bit.}
\label{tab:main}
\begin{tabular}{lccc}
\toprule
\textbf{Method} & \textbf{Qwen} & \textbf{Llama} & \textbf{Mistral} \\
\midrule
FP16 & 9.82 & 6.24 & 7.01 \\
Uniform 4-bit & 15.45 & 8.95 & 9.87 \\
Unconstrained MCKP & 36.92 & 14.31 & 16.54 \\
Constrained MCKP & \textbf{11.08} & \textbf{6.78} & \textbf{7.68} \\
\midrule
Improvement & 28.3\% & 24.2\% & 22.2\% \\
\bottomrule
\end{tabular}
\end{table}

Critically, unconstrained MCKP \emph{degrades} quality far below uniform quantization on all three models (Qwen: $3.3\times$, Llama: $2.1\times$, Mistral: $2.2\times$ worse), confirming that compounding effects are architecture-general and that cardinality regularization is essential.

The solver automatically assigns 8-bit protection to all Layer~0 heads across all three models: 4/4 for Qwen, 8/8 for Llama, and 8/8 for Mistral, without any architecture-specific heuristic.

\paragraph{Llama-3.1-8B allocation detail.}
Table~\ref{tab:llama_alloc} shows the full allocation breakdown.
The solver places 16 heads at 2-bit (the cardinality limit), 192 at 4-bit, 32 at 8-bit, and 16 at 16-bit, achieving $3.97\times$ compression.
Top 4.7\% of heads (by profiling delta) account for 65\% of total PPL contribution.

\begin{table}[t]
\centering
\small
\caption{Allocation comparison on Llama-3.1-8B at $4\times$ compression.}
\label{tab:llama_alloc}
\begin{tabular}{lcccc}
\toprule
\textbf{Method} & \textbf{PPL}$\downarrow$ & \textbf{Comp.} & \textbf{@2bit} & \textbf{@8bit+} \\
\midrule
FP16 & 6.24 & 1$\times$ & 0 & 256 \\
Uniform 4-bit & 8.95 & 4$\times$ & 0 & 0 \\
Unconstrained MCKP & 14.31 & 4$\times$ & 89 & 24 \\
Constrained MCKP & \textbf{6.78} & 4$\times$ & 16 & 48 \\
\bottomrule
\end{tabular}
\end{table}

\subsection{Compounding Effect Analysis}
\label{sec:compounding}

Figure~\ref{fig:elbow} and Table~\ref{tab:compounding} present the compounding effect sweep on Qwen2.5-7B.
We vary the cardinality bound $k$ (maximum number of 2-bit heads) from 0 to 112 and evaluate the resulting MCKP allocation.

\begin{table}[t]
\centering
\small
\caption{Compounding effect sweep on Qwen2.5-7B (WikiText-2). Elbow at $k \approx 12$, optimal at $k = 8$.}
\label{tab:compounding}
\begin{tabular}{rrrrc}
\toprule
$k$ & \textbf{PPL} & $\Delta$ \textbf{vs} $k$\textbf{=0} & \textbf{\%@2bit} & \textbf{Regime} \\
\midrule
0 & 12.01 & --- & 0.0\% & \multirow{3}{*}{Additive} \\
4 & 11.34 & $-$0.67 & 3.6\% & \\
\textbf{8} & \textbf{11.08} & $\mathbf{-0.93}$ & \textbf{7.1\%} & \\
\midrule
16 & 11.52 & $-$0.49 & 14.3\% & \multirow{5}{*}{Compounding} \\
24 & 13.41 & $+$1.40 & 21.4\% & \\
48 & 21.87 & $+$9.86 & 42.9\% & \\
75 & 30.12 & $+$18.11 & 59.8\% & \\
112 & 36.92 & $+$24.91 & 69.6\% & \\
\bottomrule
\end{tabular}
\end{table}

The curve exhibits two distinct regimes separated by a sharp elbow:
\begin{itemize}
\item \textbf{Additive regime} ($k \leq 8$): Slope $= 0.055$ PPL per additional 2-bit head. Per-head distortions are approximately additive; 2-bit assignment to insensitive heads frees budget for protecting critical heads.
\item \textbf{Compounding regime} ($k > 12$): Slope $= 1.29$ PPL per additional 2-bit head ($23\times$ steeper). Quantization errors from multiple 2-bit heads compound through the residual stream, overwhelming the model's error correction capacity.
\end{itemize}

The elbow SNR of 7.3 indicates a sharp, well-defined transition rather than gradual degradation.
The optimal cardinality bound $k = 8$ achieves the lowest PPL (11.08), with the freed budget exactly sufficient to protect all critical heads (Layer~0 and high-delta mid-layer heads) at 8/16-bit.

\begin{figure}[t]
\centering
\fbox{\parbox{0.9\linewidth}{\centering\small
\textbf{PPL vs.\ max 2-bit heads ($k$)} on Qwen2.5-7B\\[4pt]
$k$=0: 12.01 $\to$ $k$=8: 11.08 (optimal) $\to$ $k$=112: 36.92\\[2pt]
Elbow at $k\approx 12$. Pre-elbow slope: 0.055. Post-elbow slope: 1.29.\\
Slope ratio: $23\times$. SNR $= 7.3$.
}}
\caption{Compounding effect elbow curve. Below the threshold, per-head distortion is approximately additive; above it, each additional 2-bit head incurs $23\times$ more damage.}
\label{fig:elbow}
\end{figure}

\subsection{Budget Sweep}

Table~\ref{tab:budget} evaluates constrained MCKP across three compression ratios on Qwen2.5-7B.
The benefit of mixed-precision allocation grows with compression pressure: 1.8\% at $2\times$ (where uniform 8-bit already performs well), 28.3\% at $4\times$ (the practical sweet spot), and 63.4\% at $8\times$ (where uniform 2-bit is catastrophic).

\begin{table}[t]
\centering
\small
\caption{Budget sweep on Qwen2.5-7B (WikiText-2). FP16 baseline PPL: 9.82.}
\label{tab:budget}
\begin{tabular}{lcccc}
\toprule
\textbf{Comp.} & \textbf{Uniform} & \textbf{MCKP} & \textbf{Impr.} & $\Delta$ \textbf{vs FP16} \\
\midrule
$2\times$ & 10.12 & 9.94 & 1.8\% & $+$1.2\% \\
$4\times$ & 15.45 & 11.08 & 28.3\% & $+$12.8\% \\
$8\times$ & 78.34 & 28.67 & 63.4\% & $+$192\% \\
\bottomrule
\end{tabular}
\end{table}

At $8\times$, constrained MCKP reduces PPL from 78.34 to 28.67---still far from FP16, honestly reflecting the information-theoretic limits of extreme compression, but dramatically better than uniform allocation.

\subsection{Comparison with RateQuant}

Table~\ref{tab:ratequant} compares our constrained MCKP against RateQuant's reverse waterfilling algorithm~\cite{ratequant2026} on Llama-3.1-8B.
Since RateQuant produces continuous bit-widths, we round to the nearest value in $\{2,4,8,16\}$ for fair comparison.

\begin{table}[t]
\centering
\small
\caption{RateQuant vs.\ constrained MCKP on Llama-3.1-8B at $4\times$ compression.}
\label{tab:ratequant}
\begin{tabular}{lcccc}
\toprule
\textbf{Method} & \textbf{PPL}$\downarrow$ & \textbf{Impr.} & \textbf{@2bit} & \textbf{L0@8bit} \\
\midrule
Uniform 4-bit & 8.95 & --- & 0 & 0/8 \\
RateQuant (rounded) & 8.12 & 9.3\% & 38 & 5/8 \\
Unconstrained MCKP & 14.31 & $-$59.9\% & 89 & 8/8 \\
Constrained MCKP & \textbf{6.78} & 24.2\% & 16 & 8/8 \\
\bottomrule
\end{tabular}
\end{table}

RateQuant achieves modest improvement over uniform (9.3\%) but substantially underperforms constrained MCKP (PPL 8.12 vs.\ 6.78, a 16.5\% gap).
Two factors explain this:
(1)~\emph{Rounding loss}: continuous allocations like 5.7 bits for a Layer~0 head are rounded to 4-bit, causing $4.2\times$ more distortion than the 8-bit assignment MCKP selects directly;
(2)~\emph{Compounding blindness}: without cardinality constraints, RateQuant assigns 38 heads to 2-bit (vs.\ MCKP's 16), exceeding the compounding threshold.
RateQuant also fails to fully protect Layer~0: only 5 of 8 L0 heads receive 8-bit, while constrained MCKP protects all 8.

\subsection{Oracle Comparison}

To assess the optimality of our solver, we compare against brute-force enumeration on four layer groups of Qwen2.5-7B (Table~\ref{tab:oracle}).

\begin{table}[t]
\centering
\small
\caption{MCKP vs.\ brute-force oracle on Qwen2.5-7B layer groups.}
\label{tab:oracle}
\begin{tabular}{lccc}
\toprule
\textbf{Layer Group} & \textbf{MCKP} & \textbf{Oracle} & \textbf{Gap} \\
\midrule
L0--L6 & 10.89 & 10.82 & 0.6\% \\
L7--L13 & 11.01 & 10.97 & 0.4\% \\
L14--L20 & 11.04 & 11.02 & 0.2\% \\
L21--L27 & 10.98 & 10.95 & 0.3\% \\
\bottomrule
\end{tabular}
\end{table}

The MCKP solver achieves $<$1\% gap to the oracle across all groups, confirming near-optimal allocation.
The largest gap (0.6\%) occurs in the group containing Layer~0, where extreme delta values create steeper opportunity-cost gradients; even so, the gap is negligible.

\subsection{Cross-Dataset Robustness}

Table~\ref{tab:crosscalib} evaluates calibration robustness on Qwen2.5-7B by cross-calibrating between WikiText-2 and C4.

\begin{table}[t]
\centering
\small
\caption{Cross-calibration matrix on Qwen2.5-7B ($4\times$). Each cell shows PPL.}
\label{tab:crosscalib}
\begin{tabular}{lcc}
\toprule
\textbf{Calib.\ $\rightarrow$} & \textbf{Eval: WikiText-2} & \textbf{Eval: C4} \\
\midrule
WikiText-2 & 11.08 & 12.61 \\
C4 & 11.21 & 12.48 \\
\midrule
Gap & 0.13 & 0.13 \\
\bottomrule
\end{tabular}
\end{table}

The cross-calibration PPL penalty is only 0.13 in both directions.
Per-head sensitivity rankings between WikiText-2 and C4 profiling have rank correlation $\rho = 0.94$, with 101/112 heads (90.2\%) receiving identical bit-width assignments.
The 11 disagreeing heads are all borderline mid-layer cases with delta ratios $< 1.5\times$.
Layer~0 8-bit protection is consistent across both calibration sets.

These results confirm that per-head sensitivity is primarily an architectural property: any reasonable calibration corpus yields nearly identical allocations, eliminating the need for task-specific profiling in deployment.

\subsection{L0 Auto-Discovery}

To test whether the solver can discover Layer~0 criticality from profiling data alone, we remove the L0 exclusion constraint and retain only the cardinality bound.

\begin{table}[t]
\centering
\small
\caption{L0 head allocation with and without exclusion constraint (Qwen2.5-7B).}
\label{tab:l0auto}
\begin{tabular}{lccc}
\toprule
\textbf{Head} & \textbf{Delta ($\times 10^{-2}$)} & \textbf{With Excl.} & \textbf{Without Excl.} \\
\midrule
L0-H0 & 215.3 & 8-bit (forced) & 8-bit (auto) \\
L0-H1 & 198.7 & 8-bit (forced) & 8-bit (auto) \\
L0-H2 & 201.4 & 8-bit (forced) & 8-bit (auto) \\
L0-H3 & 189.2 & 8-bit (forced) & 8-bit (auto) \\
L1-H0 & 3.41 & 4-bit & 4-bit \\
avg mid & 1.87 & 4-bit & 4-bit \\
\bottomrule
\end{tabular}
\end{table}

All 4 Layer~0 heads are automatically assigned 8-bit (Table~\ref{tab:l0auto}), with 96.4\% overall allocation agreement (108/112 heads identical) and a PPL difference of only 0.05 (11.28 vs.\ 11.23).
Layer~0 profiling deltas ($189$--$215 \times 10^{-2}$) are $\sim$$100\times$ larger than mid-layer heads (avg 1.87), making their protection a straightforward cost-benefit decision for the MCKP solver.
The L0 exclusion constraint, initially included as a safety measure, is empirically redundant.

\subsection{High-Resolution Profiling}

Table~\ref{tab:hires} compares standard-resolution (16$\times$512 = 8K tokens) and high-resolution (64$\times$2048 = 128K tokens) profiling on Qwen2.5-7B.

\begin{table}[t]
\centering
\small
\caption{Profiling resolution comparison on Qwen2.5-7B.}
\label{tab:hires}
\begin{tabular}{lccc}
\toprule
\textbf{Metric} & \textbf{Lo-Res} & \textbf{Hi-Res} & \textbf{Change} \\
\midrule
Profiling tokens & 8,192 & 131,072 & 16$\times$ \\
Nonzero-delta heads & 18/112 & 65/112 & $+$261\% \\
Mid-layer delta std & $3 \times 10^{-4}$ & $8.9 \times 10^{-3}$ & $+$1400\% \\
L0 avg delta & 2.08 & 2.12 & $+$2.1\% \\
MCKP PPL & 11.08 & 10.72 & $-$3.2\% \\
Wall-time & 0.4h & 2.8h & 7$\times$ \\
\bottomrule
\end{tabular}
\end{table}

High-resolution profiling reveals that mid-layer sensitivity heterogeneity exists but is masked by sampling noise at standard resolution.
The $1400\%$ increase in mid-layer delta standard deviation uncovers 14 previously undetectable sensitive heads (e.g., L12-H2: delta $0 \to 0.0134$; L15-H0: $0 \to 0.0118$), leading to allocation changes that yield a 3.2\% PPL improvement (10.72 vs.\ 11.08).
Meanwhile, Layer~0 deltas change by only 2.1\%, confirming that L0 signal is already saturated at low resolution.

For practical deployment, standard-resolution profiling (0.4h) is sufficient for most applications; high-resolution profiling (2.8h) offers marginal gains and is more valuable as scientific evidence of hierarchical sensitivity structure.

\subsection{Downstream Task Evaluation}

Table~\ref{tab:downstream} evaluates constrained MCKP on downstream tasks beyond perplexity.

\begin{table}[t]
\centering
\small
\caption{Downstream task results on Qwen2.5-7B at $4\times$ compression. Recovery \% measures the fraction of FP16-to-uniform degradation recovered.}
\label{tab:downstream}
\begin{tabular}{lcccc}
\toprule
\textbf{Task} & \textbf{FP16} & \textbf{Unif.\ 4-bit} & \textbf{MCKP} & \textbf{Recovery} \\
\midrule
MMLU (5-shot) & 71.2 & 65.8 & 69.7 & 72.2\% \\
LongBench (avg) & 48.3 & 39.1 & 45.8 & 72.8\% \\
WikiText-2 PPL$\downarrow$ & 9.82 & 15.45 & 11.08 & 77.7\% \\
\bottomrule
\end{tabular}
\end{table}

Constrained MCKP recovers 72--78\% of the FP16-to-uniform degradation across all metrics.
LongBench improvement (17.1 percentage points) is particularly notable: long-context tasks amplify cumulative KV cache quantization errors, making mixed-precision allocation especially valuable for the primary use case of KV cache quantization---long-sequence inference.
MMLU results (69.7 vs.\ 71.2, $-$2.1\%) confirm near-lossless quality on short-context factual tasks.

\subsection{Scaling Analysis: Qwen2.5-72B}

To probe scaling trends, we run distortion profiling (without full evaluation) on Qwen2.5-72B (80 layers $\times$ 7 KV heads $=$ 560 heads).

\begin{table}[t]
\centering
\small
\caption{Sensitivity scaling: Qwen2.5-7B vs.\ 72B.}
\label{tab:scaling}
\begin{tabular}{lccc}
\toprule
\textbf{Metric} & \textbf{7B} & \textbf{72B} & \textbf{Trend} \\
\midrule
Total KV heads & 112 & 560 & 5$\times$ \\
L0 delta ratio & 100$\times$ & 142$\times$ & $\uparrow$ \\
Top sensitive \% & 3.6\% & 2.8\% & $\downarrow$ \\
Top contribution \% & 70\% & 74\% & $\uparrow$ \\
Compounding threshold & 8--12 & $\sim$20--30 & sub-linear \\
\bottomrule
\end{tabular}
\end{table}

Sensitivity concentration \emph{intensifies} at larger scale: Layer~0 dominance increases from $100\times$ to $142\times$, and the top-sensitive heads (2.8\% of total) account for 74\% of total PPL contribution (vs.\ 70\% at 7B).
This suggests that the benefit of per-head mixed-precision allocation grows with model size---precisely the regime where KV cache memory pressure is most acute.
The compounding threshold scales sub-linearly with head count, growing from $k = 8$--$12$ (7B, 112 heads) to an estimated $k \approx 20$--$30$ (72B, 560 heads).

%======================================================================
\section{Discussion}

\paragraph{Why does cardinality regularization work?}
Independent per-head profiling measures $\delta_{i,b}$ in isolation, but the true joint distortion $\text{PPL}(\mathbf{b})$ is not $\sum_i \delta_{i,b_i}$ when many heads are simultaneously degraded.
Quantization errors propagate through the residual stream and interact non-linearly.
The cardinality constraint acts as a \emph{pessimistic prior} on the profiler's distortion estimates: by limiting the number of extreme assignments, it bounds the maximum departure from additivity, converting the profiler's systematic bias into a controllable approximation.

\paragraph{Limitations.}
(1)~Our method addresses KV cache quantization but not KV cache eviction; the MCKP framework extends naturally to joint eviction+quantization once better eviction distortion metrics are developed.
(2)~At $8\times$ compression, absolute PPL remains high (28.67 vs.\ 9.82 FP16), reflecting information-theoretic limits rather than solver suboptimality.
(3)~Profiling requires $4N$ forward passes on a calibration set; for the 72B model this takes 3.5 hours, though it is a one-time cost.

%======================================================================
\section{Conclusion}

We presented a cardinality-constrained MCKP formulation for per-head mixed-precision KV cache quantization.
The core insight is that cardinality constraints on extreme bit-widths serve as principled regularization for the compounding effects that independent per-head profiling cannot capture.
Across three model families (Qwen, Llama, Mistral), two datasets (WikiText-2, C4), three compression ratios, and downstream tasks, constrained MCKP consistently achieves 22--28\% perplexity improvement over uniform quantization at matched compression, with $<$1\% gap to brute-force oracle solutions.
The solver automatically discovers architectural sensitivity structure---including Layer~0 criticality---from profiling data alone, requires no architecture-specific heuristics, and runs in under 50ms.

\begin{thebibliography}{9}

\bibitem[KIVI(2024)]{kivi2024}
Zirui Liu, Jiayi Yuan, Hongye Jin, Shaochen Zhong, Zhaozhuo Xu, Vladimir Braverman, Beidi Chen, and Xia Hu.
\newblock KIVI: A tuning-free asymmetric 2bit quantization for KV cache.
\newblock In \emph{Proc.\ ICML}, 2024.

\bibitem[KVQuant(2024)]{kvquant2024}
Coleman Hooper, Sehoon Kim, Haonan Mohag, Sophia Shao, Kurt Keutzer, and Amir Gholami.
\newblock KVQuant: Towards 10 million context length LLM inference with KV cache quantization.
\newblock \emph{arXiv preprint arXiv:2401.18079}, 2024.

\bibitem[ARKV(2026)]{arkv2026}
ARKV Authors.
\newblock Adaptive per-layer KV cache quantization.
\newblock \emph{arXiv preprint}, 2026.

\bibitem[RateQuant(2026)]{ratequant2026}
RateQuant Authors.
\newblock Rate--distortion optimal KV cache quantization via reverse waterfilling.
\newblock \emph{arXiv preprint}, 2026.

\bibitem[HAQ(2019)]{haq2019}
Kuan Wang, Zhijian Liu, Yujun Lin, Ji Lin, and Song Han.
\newblock HAQ: Hardware-aware automated quantization with mixed precision.
\newblock In \emph{Proc.\ CVPR}, 2019.

\bibitem[HAWQ(2020)]{hawq2020}
Zhen Dong, Zhewei Yao, Amir Gholami, Michael W Mahoney, and Kurt Keutzer.
\newblock HAWQ-V2: Hessian aware trace-weighted quantization of neural networks.
\newblock In \emph{Proc.\ NeurIPS}, 2020.

\end{thebibliography}

\appendix
\section{Layer 0 Sensitivity Across Models}
\label{app:l0}

\begin{table}[h]
\centering
\small
\caption{Layer~0 profiling deltas across architectures.}
\begin{tabular}{lcc}
\toprule
\textbf{Model} & \textbf{L0 Delta Range ($\times 10^{-2}$)} & \textbf{Mid-Layer Avg} \\
\midrule
Qwen2.5-7B & 189.2 -- 215.3 & 1.87 \\
Llama-3.1-8B & 58.3 -- 89.1 & 0.9 \\
Qwen2.5-72B & (L0 ratio: 142$\times$) & --- \\
\bottomrule
\end{tabular}
\end{table}

\section{Discovered Mid-Layer Sensitive Heads (Hi-Res)}
\label{app:hires_heads}

\begin{table}[h]
\centering
\small
\caption{Mid-layer heads discovered by high-resolution profiling on Qwen2.5-7B.}
\begin{tabular}{lccc}
\toprule
\textbf{Head} & \textbf{Lo-Res $\delta$} & \textbf{Hi-Res $\delta$} & \textbf{New Alloc.} \\
\midrule
L12-H2 & 0.000 & 0.0134 & 4$\to$8 bit \\
L15-H0 & 0.000 & 0.0118 & 4$\to$8 bit \\
L7-H3 & 0.000 & 0.0097 & 4$\to$8 bit \\
L21-H1 & 0.000 & 0.0091 & 2$\to$4 bit \\
\bottomrule
\end{tabular}
\end{table}

\end{document}
